% !TeX program = pdflatex
% UTF-8
\documentclass[11pt]{article}

\usepackage{iftex}
\ifPDFTeX
  \usepackage[T1]{fontenc}
\fi
\usepackage[brazilian]{babel}
\usepackage[a4paper,margin=24mm]{geometry}
\usepackage{array}
\usepackage{longtable}
\usepackage[publicar]{coop-writing}

\newcommand{\pkg}{\texttt{coop-writing}}
\newcommand{\cmd}[1]{\texttt{\textbackslash#1}}
\newcommand{\opt}[1]{\texttt{#1}}

\title{\pkg\ \coopwritingversion\\Referência Rápida (PT-BR)}
\author{Geraldo Xexéo}
\date{}

\begin{document}
\maketitle

\section*{Finalidade}

O \pkg\ é uma camada LaTeX para tratar o ciclo editorial de um documento. A informação editorial fica no código-fonte, em vez de depender da interface de edição usada. Assim, o mesmo projeto pode ser editado no Overleaf, localmente ou por Git, enquanto comentários, alterações propostas, rascunhos, anonimização, TODOs e comportamento de publicação permanecem controlados pelo LaTeX.

\section*{UTF-8 e engines}

Todos os fontes do projeto devem ser UTF-8. O LaTeX atual usa UTF-8 como codificação de entrada padrão, portanto este exemplo não carrega \texttt{inputenc}. Com pdfLaTeX ele seleciona a codificação de fonte T1; LuaLaTeX usa seu modelo Unicode nativo.

A detecção explícita de engine e a cobertura regressiva de pdfLaTeX e LuaLaTeX estão propostas na issue \#34.

\section*{Começo rápido}

Para escrita colaborativa normal:

\begin{verbatim}
\usepackage[edicao]{coop-writing}
\cwautor{alice}{blue}{Alice}
\cwautor{bruno}{red}{Bruno}
\end{verbatim}

Depois, por exemplo:

\begin{verbatim}
Esta afirmação \alice{Verifique a evidência.} precisa de revisão.

\alice[Esta frase]{Simplifique a redação.}

\brunosug{Esta frase é uma inserção proposta.}

\brunoswap{redação antiga}{redação nova}
\end{verbatim}

\section*{Modos editoriais}

Use um modo principal por compilação.

\begin{longtable}{>{\ttfamily}p{0.28\linewidth}p{0.64\linewidth}}
edicao &
Modo de trabalho. Mostra comentários, sugestões, TODOs, assuntos e rascunhos. Material que seria anonimizado é marcado, e não ocultado.\\
submeter &
Modo de submissão. Oculta as marcações editoriais e, por padrão, anonimiza o material protegido.\\
publicar &
Modo limpo de publicação. Oculta marcações, não anonimiza e mantém o lado antigo de sugestões pendentes.\\
publicaraceitando &
Modo limpo de publicação que aceita globalmente sugestões e substituições.\\
\end{longtable}

Use \opt{naoanonimizar} junto com \opt{submeter} quando a submissão não for cega. Também existem opções de nível mais baixo: \opt{comentarios}, \opt{sugestoes}, \opt{assuntos}, \opt{rascunhos}, \opt{todos}, \opt{anonimizar} e \opt{changecolor}. Para projetos normais, prefira um modo editorial principal.

Uma futura interface \texttt{chave=valor} está proposta nas issues \#32, \#33 e \#46.

\section*{Colaboradores}

\begin{verbatim}
\cwautor{alice}{blue}{Alice}
\end{verbatim}

Para cada base \texttt{nome}, o \pkg\ cria:

\begin{longtable}{>{\ttfamily\raggedright\arraybackslash}p{0.40\linewidth}p{0.52\linewidth}}
\textbackslash nome\{comentário\} & Comentário sem trecho selecionado.\\
\textbackslash nome[trecho]\{comentário\} & Comentário associado a trecho destacado.\\
\textbackslash nome r[trecho]\{comentário\}\{rótulo\} & Comentário rotulado.\\
\textbackslash nome x[...] & Comentário riscado.\\
\textbackslash nome rx[...] & Comentário rotulado e riscado.\\
\textbackslash nome sug[comentário]\{novo\} & Inserção proposta.\\
\textbackslash nome rem[comentário]\{antigo\} & Remoção proposta.\\
\textbackslash nome swap[comentário]\{antigo\}\{novo\} & Substituição proposta.\\
\textbackslash nome troca[comentário]\{antigo\}\{novo\} & Sinônimo em português de \texttt{swap}.\\
\end{longtable}

Duas famílias já existem: \cmd{cwauthor} e \cmd{cweditor}.

Use \cmd{cwsetcommwarn}\{símbolo\} para colocar um símbolo antes das marcas de comentário.

\section*{TODOs}

O default desejado é o comentário editorial curto:

\begin{verbatim}
\todo{Verificar este argumento.}
\end{verbatim}

A caixa grande é explicitamente opt-in:

\begin{verbatim}
\todo[inline]{Reescrever este parágrafo antes da submissão.}
\end{verbatim}

A cor de fundo pode ser alterada com \cmd{cwdefinetodocolor}. A issue \#45 acompanha o teste regressivo e a limpeza da implementação dessa distinção.

\section*{Citação necessária}

\begin{verbatim}
Esta afirmação quantitativa precisa de fonte.\favorcitar

Outra afirmação.\favorcitar[preferir uma revisão sistemática recente]
\end{verbatim}

\section*{Blocos de rascunho}

Na v1.8, use o nome canônico do ambiente:

\begin{verbatim}
\begin{cwdraft}[Explicação alternativa]
Este parágrafo é deliberadamente provisório.
\end{cwdraft}

\cwsetdraftcolor{cyan}
\end{verbatim}

O material de rascunho aparece durante a edição e é suprimido nos modos limpos.

\section*{Anonimização}

\begin{verbatim}
\cwanon{Universidade Federal de Exemplo}
\cwanoncite{chave}
\cwanoncitep[cap.~2]{chave}
\cwanoncitet{chave}
\end{verbatim}

Personalize as substituições com:

\begin{verbatim}
\cwsetanoncolor{violet}
\cwdefanontext{Instituição anônima}
\cwdefanoncitetext{(Anônimo, Ano)}
\cwdefanonciteptext{(Anônimo, Ano)}
\cwdefanoncitettext{Anônimo (Ano)}
\end{verbatim}

Para material nominal repetido:

\begin{verbatim}
\cwblind{\minhainstituicao}
        {Universidade Federal de Exemplo}
        {Instituição anônima}
\end{verbatim}

\section*{Assuntos e ideias principais}

\begin{verbatim}
\cwassunto{Motivação}
Este parágrafo explica por que o trabalho é necessário.

\cwmain{O pacote separa estado editorial do conteúdo do documento.}
O restante do parágrafo sustenta essa ideia principal.

\cwsetsubjectcolor{cyan}
\cwsetmaincolor{yellow}
\end{verbatim}

A apresentação de \cmd{cwmain} pode ser personalizada redefinindo \cmd{cwmainemphasis}.

\section*{Arquivos dependentes do modo}

\begin{verbatim}
\cwinput{versao-edicao}{versao-submissao}{versao-publicacao}
\cwinclude{capitulo-edicao}{capitulo-submissao}{capitulo-publicacao}

\cwinputediting{somente-edicao}
\cwinputsubmit{somente-submissao}
\cwinputpublish{somente-publicacao}

\cwincludeediting{capitulo-somente-edicao}
\cwincludesubmit{capitulo-somente-submissao}
\cwincludepublish{capitulo-somente-publicacao}
\end{verbatim}

\section*{Listas}

\begin{verbatim}
\listofcomments
\listofcitationneeds
\listofsubjects
\listoftodos

\cwreport[mode=summary]
\cwreport[type=todo,status=open]
\end{verbatim}

O relatório unificado pode filtrar por autor, tipo, severidade e estado.

\section*{Auxiliares de cor}

No modo de edição:

\begin{verbatim}
\cwcolor{blue}
Este texto fica azul.
\cwrestorecolor
\end{verbatim}

O auxiliar de baixo nível \cmd{cwsavecolor} também existe. A issue \#54 propõe tornar todas as cores do pacote configuráveis por uma única interface de usuário.

\section*{Aliases em português}

Os aliases de compatibilidade incluem \cmd{cwautor}, \cmd{cwassunto}, \cmd{favorcitar} e \cmd{listofassunto}, além das opções em português. O ambiente \texttt{rascunho} é mantido como alias de compatibilidade; a issue \#40 trata da modernização mais ampla da internacionalização.

\section*{API moderna da v1.8}

A interface preferencial usa configuração explícita por chave/valor:

\begin{verbatim}
\usepackage[mode=editing,bookmarks=false]{coop-writing}
\cwsetup{
  layout=margin-footnote,
  theme=color,
  log=summary,
  view=orientador,
  document/status=draft
}
\end{verbatim}

Itens semânticos carregam metadados e IDs estáveis:

\begin{verbatim}
\cwitem[id=metodo-1,type=comment,author=alice,
        severity=warning]{Explicar a estratégia de amostragem.}

\cwchange[id=mudanca-metodo,author=alice]
  {redação antiga}{redação nova}

\cwaccept{mudanca-metodo}
\cwresolve{metodo-1}
\cwreopen{metodo-1}

\cwreply{metodo-1}[author=bob]{Incluí a explicação.}
\end{verbatim}

Tipos e estilos reutilizáveis são independentes:

\begin{verbatim}
\cwdefinetype[command=clareza]{clarity}{
  label={Clareza}, severity=warning, color=orange
}
\cwdefinestyle{urgente}{
  severity=critical, color=red, layout=inline
}
\clareza{Explicar esta transição.}
\cwitem[type=todo,style=urgente]{Verificar o resultado.}
\end{verbatim}

Revisões longas, placeholders, visões nomeadas, metadados do documento,
resposta a revisores, errata e exportação estruturada também fazem parte da API:

\begin{verbatim}
\begin{cwrevision}[id=rev-arquitetura,author=alice]
  São permitidos parágrafos, listas e matemática display.
\end{cwrevision}

\cwplaceholder[id=fig-1,subtype=figure]{Diagrama da arquitetura}

\begin{cwview}{orientador}
Texto interno para o orientador.
\end{cwview}

\cwsetup{
  document/status=submitted,
  document/submitted-to={Revista de Exemplo},
  document/restricted=false
}

\cwloadrevisions{manuscrito}
\cwrevisionref{manuscrito}{mudanca-metodo}

\cwerratum[id=err-01,date=2026-09-24,version=1.8]
  {texto publicado}{texto corrigido}
\cwerratareport

\cwexport{itens-editoriais.tsv}
\end{verbatim}

\section*{Garantias de compatibilidade e release}

\begin{itemize}
\item \texttt{tocloft} não é mais carregado; as listas editoriais são privadas do coop-writing.
\item \texttt{hyperref} nunca é carregado automaticamente; se já estiver presente, a integração é segura.
\item pdfLaTeX e LuaLaTeX são cobertos por testes regressivos.
\item Classes padrão, memoir, KOMA-Script e fontes canônicas UFRJ/COPPE/Poli são testadas em CI.
\item \texttt{coop-writing.dtx} e \texttt{coop-writing.ins} são canônicos; style, distribuição e CTAN são gerados.
\end{itemize}

\section*{Fluxo recomendado}

Use \opt{edicao} durante a escrita e \opt{submeter} para submissão cega quando necessário. Depois das decisões editoriais, use \opt{publicar}; use \opt{publicaraceitando} quando todas as alterações propostas devam ser aceitas. Mantenha o fonte como registro autoritativo do processo editorial.

\end{document}
